\begin{ouverture}[geometrie]{15}{Symétries et demi-tour}{Partie 3 · Espace et géométrie}
{Axe · Centre · Construction · Propriétés conservées}
\objectif{Réactiver les constructions par symétrie axiale.}
\objectif{Définir et construire une symétrie centrale.}
\objectif{Utiliser les propriétés d’un demi-tour.}
\prerequis{Milieu, perpendicularité et quadrillage.}
\end{ouverture}
\pagestyle{otMain}\essentiel\fontsize{11}{15}\selectfont
\phantomsection\label{cours-15-1}
\essSection{1}{Revoir la symétrie axiale}
\begin{definition}A et A′ sont symétriques par rapport à une droite $(d)$ si $(d)$ est la médiatrice de [AA′]. Un point de l’axe est sa propre image.\end{definition}
\label{fig-cours-15-1-cours-axiale}\figureGeo{cours-axiale}

\begin{savoirfaire}[Construire sur feuille blanche]Tracer la perpendiculaire à l’axe passant par A. Appeler H son pied. Reporter HA de l’autre côté de l’axe pour placer A′. Sur un quadrillage, compter perpendiculairement à l’axe ; pour un axe diagonal, utiliser les carreaux avec soin.\end{savoirfaire}
\begin{propriete}La symétrie axiale conserve longueurs, angles, alignement et aires. L’image d’un segment s’obtient en reliant les images de ses extrémités.\end{propriete}
\clearpage
\phantomsection\label{cours-15-2}
\essSection{2}{Définir la symétrie centrale}
\begin{definition}Le symétrique A′ de A par rapport à O est le point tel que O soit le milieu de [AA′]. La symétrie centrale est un demi-tour, soit une rotation de $180^\circ$. Le centre O est sa propre image.\end{definition}
\label{fig-cours-15-2-cours-centrale}\figureGeo{cours-centrale}

\begin{propriete}Un demi-tour conserve longueurs, angles, alignement et aires. L’image d’une droite est une droite parallèle, confondue avec elle si elle passe par le centre. L’image d’un cercle est un cercle de même rayon. Ces propriétés, observées sur des figures, sont admises.\end{propriete}
\begin{attention}Ne pas confondre axe et centre. Pour un demi-tour, A, O, A′ sont alignés avec $OA=OA′$ ; pour une symétrie axiale, l’axe est perpendiculaire à [AA′]. Deux demi-tours de même centre ramènent chaque point à sa position initiale.\end{attention}
\begin{savoirfaire}[Démontrer avec un demi-tour]Deux parallèles $(d)$ et $(e)$ sont coupées par une sécante en I et J. Le demi-tour de centre M, milieu de [IJ], échange I et J et laisse la sécante en place. Il transforme $(d)$ en la parallèle à $(d)$ passant par J, c’est-à-dire $(e)$. Comme il conserve les angles, il envoie un angle alterne interne sur l’autre : ils sont égaux (chapitre 13).\end{savoirfaire}

\clearpage\rubrique{Découverte}
\begin{blocLivre}{ex-15-1}\begin{exercice}[1]{Un axe vertical}
Construire l’image du point A par rapport à l’axe vertical donné.
\label{fig-ex-15-1}\figureGeo{sym-1}
\end{exercice}\end{blocLivre}
\begin{blocLivre}{ex-15-2}\begin{exercice}[1]{Un axe horizontal}
Construire l’image de B par rapport à l’axe horizontal.
\label{fig-ex-15-2}\figureGeo{sym-2}
\end{exercice}\end{blocLivre}
\begin{blocLivre}{ex-15-3}\begin{exercice}[1]{Un demi-tour}
Construire A′, symétrique de A par rapport à O.
\label{fig-ex-15-3}\figureGeo{sym-3}
\end{exercice}\end{blocLivre}
\begin{blocLivre}{ex-15-4}\begin{exercice}[1]{Le centre lui-même}
Quelle est l’image de O par la symétrie de centre O ?
\label{fig-ex-15-4}\figureGeo{sym-4}
\end{exercice}\end{blocLivre}
\begin{blocLivre}{ex-15-5}\begin{exercice}[1]{Longueur conservée}
AB mesure 4,5 cm. Le demi-tour de centre O donne A′ et B′. Quelle est la longueur A′B′ ?
\label{fig-ex-15-5}\figureGeo{sym-5}
\end{exercice}\end{blocLivre}

\rubrique{Entraînement}
\begin{blocLivre}{ex-15-6}\begin{exercice}[2]{Axe diagonal}
Construire l’image du point $A(1;3)$ par rapport à la diagonale du quadrillage qui passe par les points $(0;0)$, $(1;1)$ et $(2;2)$.
\label{fig-ex-15-6}\figureGeo{sym-6}
\end{exercice}\end{blocLivre}
\begin{blocLivre}{ex-15-7}\begin{exercice}[2]{Image d’un triangle}
Construire l’image du triangle ABC par le demi-tour de centre O.
\label{fig-ex-15-7}\figureGeo{sym-7}
\end{exercice}\end{blocLivre}
\begin{blocLivre}{ex-15-8}\begin{exercice}[2]{Retrouver le centre}
A et A′ sont deux points images par une symétrie centrale. Construire son centre.
\label{fig-ex-15-8}\figureGeo{sym-8}
\end{exercice}\end{blocLivre}
\begin{blocLivre}{ex-15-9}\begin{exercice}[2]{Image d’une droite}
Construire l’image de (AB) par le demi-tour de centre O situé hors de (AB). Quelle relation lie les deux droites ?
\label{fig-ex-15-9}\figureGeo{sym-9}
\end{exercice}\end{blocLivre}
\begin{blocLivre}{ex-15-10}\begin{exercice}[2]{Un cercle image}
Un cercle a pour centre C et rayon 2 cm. Construire son image par la symétrie de centre O.
\label{fig-ex-15-10}\figureGeo{sym-10}
\end{exercice}\end{blocLivre}

\rubrique{Problèmes}
\begin{blocLivre}{ex-15-11}\begin{exercice}[3]{Une construction sur feuille blanche}
À partir de O et A donnés, expliquer une construction au compas de A′, image de A par le demi-tour de centre O.
\label{fig-ex-15-11}\figureGeo{sym-11}
\end{exercice}\end{blocLivre}
\begin{blocLivre}{ex-15-12}\begin{exercice}[3]{Aire et périmètre}
Un triangle de périmètre 18 cm et d’aire 12 cm$^2$ subit un demi-tour. Donner les deux mesures de son image.
\label{fig-ex-15-12}\figureGeo{sym-12}
\end{exercice}\end{blocLivre}
\begin{blocLivre}{ex-15-13}\begin{exercice}[3]{Centre d’un rectangle}
Construire le centre de symétrie d’un rectangle ABCD et expliquer quelles paires de sommets s’échangent.
\label{fig-ex-15-13}\figureGeo{sym-13}
\end{exercice}\end{blocLivre}
\begin{blocLivre}{ex-15-14}\begin{exercice}[3]{Deux demi-tours}
On applique deux fois la symétrie de centre O à un triangle. Quel triangle obtient-on ? Justifier pour un sommet.
\label{fig-ex-15-14}\figureGeo{sym-14}
\end{exercice}\end{blocLivre}
\begin{blocLivre}{ex-15-15}\begin{exercice}[3]{Un centre commun à vérifier}
A devient A′ et B devient B′ sur le dessin. Peut-il s’agir d’une même symétrie centrale ? Placer les milieux de [AA′] et de [BB′] sur le quadrillage, puis lire leurs coordonnées.
\label{fig-ex-15-15}\figureGeo{sym-15}
\end{exercice}\end{blocLivre}

\clearpage
